\documentclass[10pt,a4paper]{amsart}
\usepackage[margin=19mm]{geometry}
\usepackage{amsmath,amssymb,mathtools}
\usepackage[scr=rsfs,cal=euler]{mathalfa}
\usepackage{xfrac}
\usepackage[unicode,hidelinks,bookmarks=false]{hyperref}
\newcommand{\ie}{i.e.\ }
\newcommand{\cf}{cf.\ }
\newcommand{\resp}{respectively\ }
\newcommand{\Subsection}{\subsection}
\providecommand{\nobreakdash}{\nobreak\hbox{-}}
\setlength{\emergencystretch}{2em}
\multlinegap0pt
\usepackage{amssymb}
\usepackage[shortlabels]{enumitem}
\usepackage{xcolor}
\usepackage{float}
\newtheorem{lemma}{Lemma}
\newtheorem{definition}{Definition}
\newtheorem{proposition}{Proposition}
\newtheorem{corollary}{Corollary}
\newtheorem{theorem}{Theorem}
\newcommand{\mbb}{\mathbb}
\newcommand{\mc}{\mathcal}
\newcommand{\on}{\operatorname}
\newcommand{\wh}{\widehat}
\title{Weil-Moore anima}
\author{Dustin Clausen}
\date{Source: arXiv:2605.11950v2 (CC BY 4.0)}
\begin{document}
\maketitle

\noindent\emph{Source and licence.} Weil-Moore anima. Version arXiv:2605.11950v2 (CC BY 4.0), distributed by arXiv under the Creative Commons Attribution 4.0 International licence, http://creativecommons.org/licenses/by/4.0/, as recorded in the arXiv licence metadata for this exact version and shown on the abstract page https://arxiv.org/abs/2605.11950v2. Full licence notice: \texttt{licenses/arxiv-2605.11950-CC-BY-4.0.txt}.\par\smallskip

\noindent\textit{Editorial scope.} Counted unit: the article's theorem localtate (source0.tex lines 3158--3171). The article's complete original proof of it is delivered with it (source0.tex lines 3172--3242, one \texttt{proof} environment of 71 lines). No mathematics is written, repaired or re-derived: the statement and the proof are the article's own lines, in the article's own notation, and no retained line is substituted or re-flowed.  Retained uncounted as the source-local context the proof needs: lines 1311--1334; lines 1364--1372; lines 1736--1747; lines 1967--1983; lines 2010--2028; lines 2051--2059; lines 2207--2226; lines 2258--2267; lines 2371--2379; lines 2627--2633; lines 2681--2688; lines 3062--3078; lines 3391--3397.  Nothing was omitted from any retained span, so the package carries no omission record.  No retained line contains a citation, so the wrapper carries no bibliography and no bibliography entry.  No retained line contains a figure environment, an image-inclusion command or drawing code, so no image is delivered and the package declares no asset dependency.  Editorial formatting only, in addition: an AMS \texttt{amsart} preamble that restates the article's own declarations and macros used by the retained lines, namely the environments \texttt{lemma}, \texttt{definition}, \texttt{proposition}, \texttt{corollary}, \texttt{theorem} on separate counters, together with the standard packages those lines need.  The retained environments print in this document as Lemma 1, Definition 1, Proposition 1, Corollary 1, Theorem 1 in order of appearance, whereas the article prints the same items with the numbering of its own full document.  The complete native source of the article as distributed in the arXiv e-print for this exact version is bundled unchanged as \texttt{sources/200393/source0.tex}.
\begin{lemma}\label{transferlemma}
Let $\mc{C}$ be a Galois category, and let $\mc{D}$ be a semi-additive $\infty$-category which admits $G$-orbit and $G$-fixed point objects for all finite groups $G$.  Denote by $\on{Span}(\mc{C})$ the span category ($(2,1)$-category, actually) of $\mc{C}$, where the objects are the objects of $\mc{C}$, the maps from $X$ to $Y$ are the spans
$$X \leftarrow M \to Y$$
in $\mc{C}$, and composition is given by pullback.  Note that there are obvious functors
$$\mc{C} \to \on{Span}(\mc{C})$$
and
$$\mc{C}^{op} \to \on{Span}(\mc{C}),$$
sending each object to itself and sending a map $X\to Y$ to the span $X=X\to Y$ (resp.\ $Y\leftarrow X=X$). Then:

\begin{enumerate}
\item Let $\on{Fun}^{coshf}(\on{Span}(\mc{C}),\mc{D})$ denote the full subcategory of functors $F:\on{Span}(\mc{C})\to \mc{D}$ such that the restriction along $\mc{C}\to \on{Span}(\mc{C})$ is a cosheaf.  Then we have
$$\on{Fun}^{coshf}(\on{Span}(\mc{C}),\mc{D})\overset{\sim}{\rightarrow}\on{coSh}(\mc{C};\mc{D})$$
via restriction along $\mc{C} \to \on{Span}(\mc{C})$.
\item Let $\on{Fun}^{shf}(\on{Span}(\mc{C}),\mc{D})$ denote the full subcategory of functors $F:\on{Span}(\mc{C})\to \mc{D}$ such that the restriction along $\mc{C}^{op}\to \on{Span}(\mc{C})$ is a sheaf.  Then we have
$$\on{Fun}^{shf}(\on{Span}(\mc{C}),\mc{D})\overset{\sim}{\rightarrow}\on{Sh}(\mc{C};\mc{D})$$
via restriction along $\mc{C}^{op} \to \on{Span}(\mc{C})$.
\item Let $F \in \on{Fun}(\on{Span}(\mc{C}),\mc{D})$.  Any two of the following three properties imply the third:
\begin{enumerate}
\item $F\in \on{Fun}^{coshf}(\on{Span}(\mc{C}),\mc{D})$;
\item $F\in \on{Fun}^{shf}(\on{Span}(\mc{C}),\mc{D})$;
\item for all connected $X\in\mc{C}$ and all subgroups $G\subset \on{Aut}(X)$, the natural ``transfer'' map $F(X)_G \to F(X)^G$, associated to the induced $G$-action on $F(X)$, is an isomorphism.
\end{enumerate}
\end{enumerate}
\end{lemma}

\begin{definition}\label{trfunctor}
Let $\mc{C}$ be a Galois category, $\mc{D}$ a semi-additive $\infty$-category which admits $G$-orbit and $G$-fixed point objects for all finite groups $G$, and let $F\in \on{coSh}(\mc{C};\mc{D})$.  Define the presheaf
$$F^{tr} \in \on{Fun}(\mc{C}^{op},\mc{D})$$
as the image of $F$ under the composition
$$\on{coSh}(\mc{C};\mc{D})\simeq\on{Fun}^{coshf}(\on{Span}(\mc{C}),\mc{D})\rightarrow \on{Fun}(\mc{C}^{op},\mc{D}),$$
where the first map is the inverse to the isomorphism of 1 of Lemma \ref{transferlemma} and the second map is composition with the natural map $\mc{C}^{op}\to\on{Span}(\mc{C})$ described in Lemma \ref{transferlemma}.

Dually, if $F$ is a sheaf we get a co-presheaf $F^{tr}$.  Note that by construction $F$ and $F^{tr}$ agree on objects; moreover the association $F\mapsto F^{tr}$ is functorial in $\mc{C}$ (for exact functors) and $\mc{D}$ (for additive functors), via composition.
\end{definition}

\begin{definition}\label{maprcf}
Let $(\mc{C},F,\alpha)$ and $(\mc{C}',F',\alpha')$ be two refined class formations (condensed or not: the discussion is the same in either case).  A morphism 
$$f:(\mc{C},F,\alpha)\to (\mc{C}',F',\alpha')$$
consists of the following data:
\begin{enumerate}
\item An exact functor $f^{-1}:\mc{C}'\to\mc{C}$;
\item A natural transformation $f_\natural: f_\ast F \to F'$;
\end{enumerate}
subject to the condition that $\alpha'\circ f_\natural$ should agree with the map induced on $f_\ast$ by $\alpha:F\to\mbb{Z}$.

Here $f_\ast G := G\circ f^{-1}$ for any functor $G$ from $\mc{C}$ to some target $\infty$-category; note in particular that $f_\ast$ sends (co)sheaves to (co)sheaves, and sends constant functors to constant functors.  (It also commutes with the $(-)^{tr}$ operation from Definition \ref{trfunctor}, since $f^{-1}$, being exact, also induces an additive functor on span categories.)
\end{definition}

\begin{lemma}\label{gerbe}
Consider the $\infty$-category $\mc{A}$ of pairs $(M,\alpha)$ where $M\in\on{D}(\on{CondAb})$ lives in degrees $[0,1]$ and $\alpha:H_0(M)\simeq\mbb{Z}$.  Then:
\begin{enumerate}
\item There is a fully faithful embedding $\mc{A} \to \on{CondAn}$, sending
$$(M,\alpha) \mapsto (\Omega^\infty M)\times_{\mbb{Z}} \{1\},$$
where the map $\Omega^\infty M\to\mbb{Z}$ is induced by $\alpha.$
\item The essential image of this functor consists of the \emph{abelian gerbes}: namely, those condensed anima isomorphic to $BA$ for some condensed abelian group $A$.
\item The inverse to this functor sends
$$X\mapsto (\tau_{\leq 1}(\mbb{Z}[X]),\epsilon);$$
for $X$ an abelian gerbe, where $\epsilon:\tau_{\leq 1}\mbb{Z}[X] \to \mbb{Z}$ is induced by $X\to \ast$.  More generally, for $X\in\on{CondAn}$ arbitrary and $(M,\alpha)\in\mc{A}$, we have that maps
$$(\tau_{\leq 1}(\mbb{Z}[X]),\epsilon) \to (M,\alpha)$$
(in the obvious sense) are in natural equivalence with maps
$$X \to (\Omega^\infty M)\times_{\mbb{Z}} \{1\}.$$
\item For any $X\in\on{CondAn}$ with $\pi_0X=\ast$, there is an initial abelian gerbe $X^{ab}$ with a map $X\to X^{ab}$, namely that given by
$$X^{ab} = (\tau_{\leq 1}\mbb{Z}[X])\times_{\mbb{Z}}\{1\}.$$
\end{enumerate}
\end{lemma}

\begin{proposition}\label{rcfgerbe}
For a Galois category $\mc{C}$, the data $(F,\alpha)$ of a condensed refined class formation on $\mc{C}$ is equivalent to the datum of a functor
$$\mc{G}:\mc{C}_{conn}\to\on{CondAn}$$
where $\mc{C}_{conn}\subset\mc{C}$ is the full subcategory of connected objects, satisfying the following properties:
\begin{enumerate}
\item $\mc{G}$ lands in the full subcategory of abelian gerbes;
\item For all $X\in\mc{C}_{conn}$ and all subgroups $G\subset\on{Aut}(X)$, the map in $\on{CondAn}$
$$\mc{G}(X)_G \to \mc{G}(X/G)$$
induces an isomorphism on abelianization
$$(\mc{G}(X)_G)^{ab} \overset{\sim}{\rightarrow} \mc{G}(X/G).$$
\item For all $X\in\mc{C}_{conn}$ and all subgroups $G\subset\on{Aut}(X)$, the transfer map of condensed abelian groups
$$H_1\mc{G}(X/G)\overset{\sim}{\rightarrow} \left(H_1\mc{G}(X)\right)^G$$
(defined by passing through Lemma \ref{gerbe} using 2) is an isomorphism.
\end{enumerate}
In other words, $\mc{G}$ is the datum of a cosheaf of abelian gerbes on $\mc{C}_{conn}$ satisfying the extra condition about the transfer map being an isomorphism.

Moreover, given a condensed refined class formation $C$ encoded in such data, and another $C'$ similarly encoded, then giving a map of condensed refined class formations $C \to C'$ (Definition \ref{maprcf}) is the same as giving an exact functor $f^{-1}:\mc{C}'\to\mc{C}$ plus a natural transformation
$$f_\natural: f_\ast\mc{G} \to \mc{G}'.$$
\end{proposition}

\begin{definition}\label{rcftocond}
Let $C=(\mc{C},F,\alpha)$ be a condensed refined class formation, and suppose $\mc{G}:\mc{C}_{conn}\to\on{CondAn}$ is the corresponding cosheaf of abelian gerbes, as in Proposition \ref{rcfgerbe}.  Define a condensed anima $\mathcal{W}_C$ by setting, for $S$ extremally disconnected profinite, $\mathcal{W}_C(S)$ to be the anima of pairs $(f^{-1},f_\natural)$ where:
\begin{enumerate}
\item $f^{-1}$ is an exact functor $\mc{C} \to \on{fEt}_S$;
\item $f_\natural$ is a natural transformation $f_\ast \tau_S \to \mc{G}$, i.e.\ a family of maps of condensed anima
$$f^{-1}(U)\to \mc{G}(U)$$
natural in $U\in\mc{C}_{conn}$.
\end{enumerate}
\end{definition}

\begin{corollary}\label{weildeterminesrcf}
The functor sending a condensed refined class formation $C = (\mc{C},F,\alpha)$ to the condensed anima $\mc{W}_C$
from Definition \ref{rcftocond} induces an equivalence between the following two $(2,1)$-categories:
\begin{enumerate}
\item The full subcategory of condensed refined class formations spanned by those $(\mc{C},F,\alpha)$ such that for all connected $U\in\mc{C}$, we have:
\begin{enumerate}
\item $H_1F(U)$ is locally compact Hausdorff;
\item $\on{ker}(H_1F(U)\to H_1F(\ast))$ is compact;
\item Every open subgroup of $H_1F(U)$ of finite index contains $\on{im}(H_1F(V)\to H_1F(U))$ for some connected $V\in\mc{C}$ mapping to $U$. (This property is the ``existence theorem'' for the underlying class formation.)
\end{enumerate}

\item The full subcategory of $\on{CondAn}$ spanned by those objects of the form $BW$ where $W$ is a locally compact Hausdorff group such that:
\begin{enumerate}
\item $W\overset{\sim}{\rightarrow}\varprojlim_N W/N^c$ where $N$ runs over all open normal subgroups of $W$ of finite index and $N^c$ is the closure of the commutator subgroup of $N$, and the transition maps in this inverse system have compact kernel;
\item For any open subgroup $H\subset W$ of finite index and any open normal subgroup $N\subset H$ of finite index, the transfer map induces an isomorphism of locally compact Hausdorff groups
$$H/H^c\overset{\sim}{\rightarrow} (N/N^c)^{H/N}.$$

\end{enumerate}
\end{enumerate}
\end{corollary}

\begin{proposition}\label{ztype}
The equivalence in Corollary \ref{weildeterminesrcf} restricts to an equivalence between the following two $(2,1)$-categories:

\begin{enumerate}
\item The category of refined class formations of $\mbb{Z}$-type;
\item The full subcategory of $\on{CondAn}$ spanned by those condensed anima of the form $BW$, where $W$ is of the following type: take a profinite group $\Gamma$ of \emph{strict cohomological dimension $2$} and a surjective homomorphism $\Gamma\twoheadrightarrow \wh{\mbb{Z}}$ and set
$$W= \Gamma\times_{\wh{\mbb{Z}}}\mbb{Z}.$$
\end{enumerate}
Moreover, in this case we have $H^i(U;\mbb{R}/\mbb{Z})=0$ for all $i\geq 2$ and all $U\in\on{fEt}_{BW}$, so that the associated condensed refined class formation is given by $F=R\Gamma(-;\mbb{R}/\mbb{Z})^\vee$ without any truncation needed.
\end{proposition}

\begin{proposition}\label{extvanishing}
Let $C$ be a refined class formation which is either of $\mbb{Z}$-type, of $\mbb{R}$-type, or of compact type, with associated Weil groupoid $\mc{W}_C\simeq BW$.  If $C$ is of $\mbb{Z}$-type, assume furthermore the following technical hypothesis: the kernel $W^1$ of the corresponding surjective homomorphism
$$W \twoheadrightarrow \mbb{Z},$$
which is a profinite group (see Proposition \ref{ztype}), has cohomological dimension $\leq 1$ with respect to torsion coefficients.


Then for any discrete $W/W^\circ$-module $M$ and any $i\geq 3$, we have
$$\on{Ext}^i_{W/W^\circ}(M,c_W)=0.$$
\end{proposition}

\begin{theorem}\label{strictcohdim}
Let $X$ be a Weil-Moore anima with underlying class formation either of $\mbb{Z}$-type, $\mbb{R}$-type, or compact type, and let $W=\pi_1X$ denote ``the" associated Weil group, with profinite completion $\wh{W}$ (the Galois group).  If $X$ is of $\mbb{Z}$-type, assume the same additional technical assumption from Proposition \ref{extvanishing} is satisfied. Then for any discrete group $M$ with $W$-action, we have:
\begin{enumerate}
\item $H^i(X;M)=0$ for $i>3$;
\item If either $M$ comes from a $\wh{W}$-action or $M$ is torsion, then $H^3(X;M)=0$.
\end{enumerate}
\end{theorem}

\begin{proposition}\label{pfinite}
Let $X$ be a Weil-Moore anima with underlying class formation either of $\mbb{Z}$-type, $\mbb{R}$-type, or compact type, and set $W=\pi_1X$.  For a prime $p$, the following properties are equivalent:
\begin{enumerate}
\item If $A$ denotes the sheaf of abelian groups on $\on{fEt}_X$ describing the class formation associated to $X$, then for all $Y\to X$ in $\on{fEt}_X$ connected, both $A(Y)[p]$ and $A(Y)/pA(Y)$ are finite abelian groups; or equivalently, the derived $A(Y)/p$ is finite in each degree.
\item For all $Y\to X$ finite etale connected, $R\Gamma(Y;\mbb{F}_p)\in\on{Perf}(\mbb{F}_p)$.
\item Let $R$ be a $p$-power torsion $E_1$-ring and let $M\in \on{Perf}(R)$ with $\Omega X$-action.  If we are in the $\mbb{Z}$-case, assume that $M$ descends to $B\wh{W}$.  Then $R\Gamma(X;M)\in \on{Perf}(R)$.
\end{enumerate}
\end{proposition}

\begin{lemma}\label{lisse}
Let $X$ be a Weil-Moore anima, with underlying class formation of $\mbb{Z}$-type, $\mbb{R}$-type, or compact type, with $W=\pi_1X$.  Then:

\begin{enumerate}
\item The global sections functor $\on{D}_{et}(X;\mbb{Z}) \to \on{D}(\on{CondAb})$ lands in $\on{D}(\mbb{Z})$ and preserves colimits.
\item For every connected finite etale $Y\to X$, the object $\mbb{Z}[h_Y]\in \on{D}_{et}(X;\mbb{Z})$ corepresenting sections over $Y$ is a compact object.  If $Y\to X$ corresponds to the open subgroup $H\subset W=\pi_1X$ of finite index, then $\mbb{Z}[h_Y]$ lives in degree $0$ and corresponds to $\mbb{Z}[W/H]$ as a $W$-module.
\item Denote by $\on{Lisse}(X)\subset\on{D}_{et}(X;\mbb{Z})$ the thick subcategory generated by the objects $\mbb{Z}[h_Y]$ as above.  Then $M\in\on{Lisse}(X)$ if and only if $H_iM=0$ for all but finitely many $i$, and $H_iM$ corresponds to a finitely generated discrete abelian group with continuous $\wh{W}$-action for all $i$.

\item The induced colimit-preserving functor
$$\on{Ind}(\on{Lisse}(X))\to \on{D}_{et}(X;\mbb{Z})$$
is fully faithful.
\item $\on{Ind}(\on{Lisse}(X))$ inherits the t-structure from $\on{D}_{et}(X;\mbb{Z})$, i.e.\ the essential image of the functor in 4 is closed under t-truncation.
\item Suppose $X$ is either compact-type or $\mbb{R}$-type.  Then for all $n\in\mbb{Z}$, the full subcategories of $t_{\leq n}$-truncated objects in $\on{Ind}(\on{Lisse}(X))$ and $\on{D}_{et}(X;\mbb{Z})$ are identified via the above fully faithful functor, and (hence) $\on{D}_{et}(X;\mbb{Z})$ is the left-completion of $\on{Ind}(\on{Lisse}(X))$.
\item Suppose $X$ is of profinite type (meaning, the class groups in the underlying class formation are profinite abelian groups, or equivalently $X\simeq BW$ where $W$ is profinite).  Then $\on{Ind}(\on{Lisse}(X))=\on{D}_{et}(X;\mbb{Z})$.
\end{enumerate}

\end{lemma}

\begin{lemma}\label{alttr}
Let $\mc{C}$ be a Galois category, $\mc{D}$ a $\mbb{Z}$-linear presentable $\infty$-category, and $\mc{F}\in\on{Sh}(\mc{C};\mc{D})$.  Then there is a natural isomorphism between the following two functors $\mc{C}\to\mc{D}$:
\begin{enumerate}
\item The functor $U\mapsto \Gamma(\mc{C};\mbb{Z}[h_U]\otimes\mc{F})$;
\item The functor $\mc{F}^{tr}$ from Definition \ref{trfunctor}.
\end{enumerate}
\end{lemma}

\begin{theorem}\label{localtate}
Let $p$ be a prime and $K$ a local field of characteristic $\neq p$.  Let $X$ be a Weil-Moore anima for $K$, with $W=\pi_1X$ the Weil group.  Then:
\begin{enumerate}
\item There is a natural map
$$[X]:R\Gamma(X;\left(\mbb{Q}_p/\mbb{Z}_p\right)(1)[2])\to \mbb{Q}_p/\mbb{Z}_p$$
which is an isomorphism in degree $0$, the bottom (homological) degree.
\item For every $p$-power torsion $E_1$-$\mbb{Z}$-algebra $R$ and every $M\in\on{Perf}(R)$ with $\Omega X$-action, assumed to come from a $\wh{W}$-action if $K$ is nonarchimedean, the $R$-$R$-bimodule map
$$R\Gamma(X;M(1)[2])\otimes_\mbb{Z} R\Gamma(X;M^\vee)\to R\Gamma(X;R(1)[2])\overset{[X]\otimes_{\mbb{Z}}R[-1]}{\longrightarrow} R$$
induces an isomorphism
$$R\Gamma(X;M^\vee)\simeq R\Gamma(X;M(1)[2])^\vee$$
in $\on{Perf}(R^{op})$, where $(-)^\vee$ is the functor of left $R$-linear maps to $R$, viewed as landing in $\on{Perf}(R^{op})$ via the right $R$-module structure on $R$.

\end{enumerate}
\end{theorem}

\begin{proof}
Consider the sheaf on $\on{fEt}_X$ given by
$$\mathcal{F}:= \left(R\Gamma(-;(\mbb{Q}_p/\mbb{Z}_p)(1)[2])^\vee\right)^{tr},$$
so that, in terms of the fundamental cosheaf $F$ from the refined class formation, we have
$$\mc{F} = (F^{tr})_{\wh{p}}(-1)[-2].$$
By construction of the class formation the (unsheafified) $H_0(F^{tr})$ is $\mbb{Z}^{tr}$ and $H_1(F^{tr})$ is $\mbb{G}_m$ and all the other $H_i$ are $0$.  Thus, $(F^{tr})_{\wh{p}}$ has top nonvanishing homology group given by $H_2 = \mbb{Z}_p(1)$.  Tate untwisting and shifting down by two, we deduce that
$$H_0\mc{F} = \mbb{Z}_p$$
and that this is the top nonvanishing homology group sheaf of $\mc{F}$.  In particular, $1\in\mbb{Z}_p$ corresponds to a canonical map

$$[X]: R\Gamma(X;(\mbb{Q}_p/\mbb{Z}_p)(1)[2])\to\mbb{Q}_p/\mbb{Z}_p$$
which is an isomorphism in bottom degree $0$ by Pontryagin duality.  This proves 1.

For 2, in the non-archimedean case let us replace $X$ by $\wh{X}$, so that in any case we are either $\mbb{R}$-type or compact type. Let
$$\Gamma^!:\on{D}(\mbb{Z})[p^\infty] \to \on{D}_{et}(X;\mbb{Z})[p^\infty]$$
be the right adjoint to the global sections functor in the opposite direction (on $p$-primary torsion discrete derived coefficient systems).  We claim that the map
$$(\mbb{Q}_p/\mbb{Z}_p)(1)[2] \to \Gamma^!(\mbb{Q}_p/\mbb{Z}_p)$$
adjoint to the map of 1 is an isomorphism, or equivalently (checking mod $p$ and twisting) that the induced map
$$\mbb{F}_p \to (\Gamma^!\mbb{F}_p)(-1)[-2]$$
is an isomorphism.

To prove this, note that $\Gamma^!\mbb{F}_p$ is bounded above in the t-structure by the cohomological dimension estimate in Lemma \ref{strictcohdim}.  Thus, by Lemma \ref{lisse} the above map is an isomorphism if and only if induces an isomorphism of sheaves on $\on{fEt}_X$ when we map in from $\mbb{F}_p[h_Y]$ for arbitrary $Y\to X$ finite etale.  When we implement this, by Lemma \ref{alttr} the object $\Gamma^!\mbb{F}_p(-1)[-2]$ gives rise to the sheaf on $\on{fEt}_X$ given by
$$Y\mapsto (R\Gamma(Y;\mbb{F}_p(1)[2])^\vee)^{tr},$$
or in other words, in terms of the fundamental cosheaf $F$ coming from the refined class formation,
$$Y\mapsto F^{tr}(Y)/p(-1)[-2],$$\
which is exactly the mod $p$ version of the $\mc{F}$ we analyzed above.  Moreover, the map $\mbb{F}_p\to (\Gamma^!\mbb{F}_p)(-1)[-2]$ is we are trying to show is an isomorphism is just selecting a global section of this sheaf on $\on{fEt}_X$, and this global section is the one we wrote down above in the proof of 1.

On the other hand, the source $\mbb{F}_p$ gives rise to the sheaf on $\on{fEt}_X$ given by
$$Y\mapsto R\Gamma(Y;\mbb{F}_p),$$
which in terms of $F$ is
$$Y\mapsto (F/p)^\vee.$$
By construction, this map of sheaves on $\on{fEt}_X$
$$R\Gamma(-;\mbb{F}_p) \to (R\Gamma(-;\mbb{F}_p(1)[2])^\vee)^{tr}$$
is an isomorphism in degree $0$ on global sections, and we want to prove it's an isomorphism.  We will do this by seeing that it's an isomorphism on stalks.

Consider the natural map
$$\mbb{F}_p \to (F/p)^\vee$$
induced by $\alpha:F\to\mbb{Z}$, whose fiber is given by the (derived) $p$-torsion in the Pontryagin dual to the class group.  By the existence theorem, when we pass to stalks we can replace the class group by its connected component.  Thus in the nonarchimedean case $\mbb{F}_p \to (F/p)^\vee$ is an isomorphism on stalks, while in the archimedean case we get that the stalks of $(F/p)^\vee$ are given by $\mbb{F}_p$ in degree $0$ and $\mbb{F}_p(-1)$ in degree $-2$.  We can likewise analyze $F^{tr}/p(-1)[-2]$.  When $F$ is nonarchimedean the sheafified $H_0F^{tr}$ is $\mbb{Q}$ so doesn't contribute mod $p$; moreover multiplication by $p$ on $H_1F=\mbb{G}_m$ is surjective on stalks because $\on{char}(F)\neq p$.  Thus we get that for the stalks of $F^{tr}/p(-1)[-2]$ there is only the $\mbb{F}_p$ in degree $0$ we already discussed at the beginning of this proof.  But when $F$ is archimedean we additionally get an $\mbb{F}_p(-1)$ in degree $-2$ (if we normalize so that $H_0F^{tr}=\mbb{Z}$, and not $\frac{1}{2}\mbb{Z}$, on the complex numbers).

In the nonarchimedean case, we are done: a map from $\mbb{F}_p$ to itself is either $0$ or an isomorphism, but our map was nonzero on global sections by construction.  Thus it must be an isomorphism.

In the archimedean case, because we are checking on stalks we reduce to the case $K=\mbb{C}$.  Then the Tate twist and the $(-)^{tr}$ operation can be ignored, and our map is simply a map
$$R\Gamma(X;\mbb{F}_p) \to R\Gamma(X;\mbb{F}_p[2])^\vee$$
induced by cup product followed by a nonzero linear form on $H^2$.  Because of the simple structure of cohomology (nonzero except in degrees $0$ and $2$, where it's one-dimensional) this forces the map to be an isomorphism, as desired.

To sum up, we have shown that the map in 1 induces an isomorphism
$$(\mbb{Q}_p/\mbb{Z}_p)(1)[2] \overset{\sim}{\rightarrow} \Gamma^!(\mbb{Q}_p/\mbb{Z}_p).$$
Next we claim that\
$$\Gamma^!:\on{D}(\mbb{Z})[p^\infty] \to \on{D}_{et}(X;\mbb{Z})[p^\infty]$$
preserves all colimits, or equivalently that the natural map
$$M\otimes_\mbb{Z} \Gamma^!(\mbb{Q}_p/\mbb{Z}_p[-1]) \to \Gamma^!(M)$$
is an isomorphism for all $M\in\on{D}(\mbb{Z})[p^\infty]$.  We can again test mod $p$ and thus reduce to showing
$$M\otimes_{\mbb{F}_p} \Gamma^!(\mbb{F}_p)\overset{\sim}{\rightarrow} \Gamma^!(M)$$
for all $M\in\on{D}(\mbb{F}_p)$.  Because $\Gamma^!(\mbb{F}_p)$ (as calculated above) is just a shift and a twist, the functor on the left commutes with the inverse limit of the Postnikov tower of $M$; the functor on the right evidently does as well because $\Gamma^!$ is a right adjoint.  Thus we can reduce to the case of $M$ bounded above, when the left-hand side is also bounded above.  Then again by Lemma \ref{lisse} we can test on maps from objects of the form $\mbb{F}_p[h_Y]$ for $Y\to X$ finite etale, in which case the claim follows from the mod $p$ finiteness property of $Y$ from Proposition \ref{pfinite}.

To sum up once more, we have shown that\
$$\Gamma^!(M) \simeq M(1)[2]$$
for all $M\in\on{D}(\mbb{Z})[p^\infty]$.  As both $\Gamma$ and $\Gamma^!$ preserve colimits, the adjunction between them passes to $R$-module objects for all $E_1$-algebras $R$ in $\on{D}(\mbb{Z})[p^\infty]$, i.e.\ all $p$-primary torsion $E_1$-$\mbb{Z}$-algebras $R$.  Now, let $M$ be as in 2: an $R$-linear coefficient system on $X$ whose fiber at $\ast \to X$ lies in $\on{Perf}(R)$.  Then we can calculate the $R^{op}$-module of $R$-linear maps
$$ M(1)[2] \to \Gamma^!(R)$$
 in two ways:
\begin{enumerate}
\item By the adjunction between $\Gamma^!$ and $\Gamma$, they are the same as $R$-linear maps
$$ R\Gamma(X;M(1)[2]) \to R,$$
so we get
$$R\Gamma(X;M(1)[2])^\vee.$$
\item By the above calculation $\Gamma^!(R)\simeq R(1)[2]$, we get
$$R\Gamma(X;M^\vee).$$
\end{enumerate}
This gives an isomorphism
$$R\Gamma(X;M^\vee)\simeq R\Gamma(X;M(1)[2])^\vee,$$
which by a simple unwinding is induced by the pairing described in 2.
\end{proof}

\end{document}
